\documentclass[10pt,a4paper]{amsart}
\usepackage[margin=19mm]{geometry}
\usepackage{amsmath,amssymb,mathtools}
\usepackage[scr=rsfs,cal=euler]{mathalfa}
\usepackage{xfrac}
\usepackage[unicode,hidelinks,bookmarks=false]{hyperref}
\newcommand{\ie}{i.e.\ }
\newcommand{\cf}{cf.\ }
\newcommand{\resp}{respectively\ }
\newcommand{\Subsection}{\subsection}
\providecommand{\nobreakdash}{\nobreak\hbox{-}}
\setlength{\emergencystretch}{2em}
\multlinegap0pt
\usepackage{bm}
\usepackage{bbm}
\usepackage{dsfont}
\usepackage{xspace}
\usepackage{cancel}
\usepackage{mathtools}
\usepackage{amsthm}
\makeatletter
\@ifundefined{theorem}{\newtheorem{theorem}{Theorem}}{}
\@ifundefined{corollary}{\newtheorem{corollary}[theorem]{Corollary}}{}
\makeatother
\title[Discrete Space-Time Wave Kernels and Trace Identities on Reg]{Discrete Space-Time Wave Kernels and Trace Identities on Regular Graphs}
\author{Amar Ba\v{s}i\'c, Lejla Smajlovi\'c, Zenan \v{S}abanac}
\date{Source: arXiv:2606.27075v1 (CC BY 4.0)}
\begin{document}
\maketitle

\noindent\emph{Source and licence.} Discrete Space-Time Wave Kernels and Trace Identities on Regular Graphs. Version arXiv:2606.27075v1 (CC BY 4.0), distributed under the Creative Commons Attribution 4.0 International licence, as recorded in the arXiv licence metadata for this exact version and shown on the abstract page https://arxiv.org/abs/2606.27075v1. Full rights notice: licenses/arxiv-2606.27075-CC-BY-4.0.txt.\par\smallskip

\noindent\textit{Editorial scope.} Counted unit: the Corollary whose source label is \texttt{cor:Kn\_c\_identities} in the article (source0.tex lines 1104--1147, source label \texttt{cor:Kn\_c\_identities}), delivered together with the article's own complete original proof of it (source0.tex lines 1149--1259, one \texttt{proof} environment of 111 lines). No mathematics is written, repaired or re-derived: statement and proof are the article's own text line for line. Required context retained uncounted: eq:edge (lines 78--83); eq. trace fla (lines 174--181); eq. bm complete graph (lines 1006--1016). Every definition, display and label of those lines is retained unchanged. Omitted from this package and not redistributed: 1--77, 84--173, 182--1005, 1017--1103, 1148--1148, 1260--1627. Nothing inside a retained excerpt is omitted or altered: every retained body is delivered exactly as the article's lines are written, so no omission receipt of any kind accompanies this package. Citations: the retained excerpts carry no citation command at all, so no bibliography entry of the article is transcribed and no reference is redistributed. Reference adaptation: none was needed and none was made; every reference inside the delivered unit resolves inside it, so no unresolved reference or citation is produced. Printed slips kept literal: no printed word, symbol, hypothesis or number of the counted statement or of its proof was altered. Layout and class: the article's own class is \texttt{amsart} with packages including graphicx, color, amsmath, amssymb, amsthm, enumitem, bbm, nag, onlyamsmath, csquotes, url; the wrapper is the lane's pinned \texttt{amsart} editorial wrapper at 10pt on A4, which restates the theorem-like environments the retained text needs and adds the article's own macro definitions that the retained text needs (none), with extra wrapper packages (bm, bbm, dsfont, xspace, cancel, mathtools). The delivered numbering is the unit's own. Assets: no retained span contains a figure environment, a graphics-inclusion command, a PDF-page inclusion, a table or a TikZ picture; no image file is copied into this package, the package declares no asset dependency and no image file is redistributed with it. Licence: version arXiv:2606.27075v1 of this article is distributed under the Creative Commons Attribution 4.0 International licence, as recorded in the arXiv licence metadata for this exact version and shown on the abstract page https://arxiv.org/abs/2606.27075v1. A full rights notice is delivered at licenses/arxiv-2606.27075-CC-BY-4.0.txt. Characters: every retained excerpt is pure ASCII, so the delivered engine, XeLaTeX with its default Latin Modern text font, reports no missing character.
\begin{equation}\label{eq:edge}
a\leq \begin{cases}
0, & \text{if } b>0 \\
-|b|\max\{\delta_{X={\textrm{finite}}}\lambda_{\max}, (\sqrt{q}+1)^2\}, & \text{if } b<0,
\end{cases}
\end{equation}

\begin{equation}\label{eq. trace fla}
\begin{split}
        \sum_{j=1}^N& 
\mu_j^k 
\psi_j(x_0)\overline{\psi_j(x)} = d_{a,b}^k \sum_{m=0}^{k} (-1)^{m} b_m(x) \, q^{-\tfrac{m}{2}} \, I_{m}^{c_{a,b}}(k)\\
&=d_{a,b}^k \left( \delta_{x_0=x}I_0^{c_{a,b}}(k) + 2 \sum_{m=1}^{k} (-1)^{m} \Bigl[T_m\left(\frac{A}{2\sqrt{q}}\right)\delta_{x_0}\Bigr](x)\, I_{m}^{c_{a,b}}(k)\right),
\end{split}
\end{equation}

\begin{equation}\label{eq. bm complete graph}
b_m(x_\ell)
=
\frac{2(n-2)^{m/2}}{n}
\left[
T_m\!\left(\frac{n-1}{2\sqrt{n-2}}\right)
+
\eta_\ell
T_m\!\left(-\frac{1}{2\sqrt{n-2}}\right)
\right],
\end{equation}

\begin{corollary}\label{cor:Kn_c_identities}
Let $n\ge3$. Then, for every $k\in\mathbb N_0$ and every real number $c$, one has
\begin{equation}\label{eq:Kn_first_identity}
\begin{split}
\sum_{m=1}^{k}
(-1)^m
&\left(
T_m\!\left(\frac{n-1}{2\sqrt{n-2}}\right)
-
T_m\!\left(-\frac{1}{2\sqrt{n-2}}\right)
\right)
I_m^{c}(k)\\
&\qquad =
\frac12
\left[
\left(1-\frac{(n-1)c}{2\sqrt{n-2}}\right)^k
-
\left(1+\frac{c}{2\sqrt{n-2}}\right)^k
\right],
\end{split}
\end{equation}
and 
\begin{equation}\label{eq:Kn_second_identity}
\begin{split}
\sum_{m=1}^{k}&
(-1)^m
\left(
T_m\!\left(\frac{n-1}{2\sqrt{n-2}}\right)
+
(n-1)T_m\!\left(-\frac{1}{2\sqrt{n-2}}\right)
\right)
I_m^{c}(k)\\
&\quad =
\frac12
\left[
\left(1-\frac{(n-1)c}{2\sqrt{n-2}}\right)^k
+
(n-1)
\left(1+\frac{c}{2\sqrt{n-2}}\right)^k
\right]
-\frac{n}{2}I_0^c(k).
\end{split}
\end{equation}
\end{corollary}

\begin{proof}
An orthonormal basis of eigenfunctions of the combinatorial Laplacian on
$K_n$ is given by
\[
\psi_j(x_\ell)=\frac1{\sqrt n}e^{2\pi i j\ell/n},
\qquad j,\ell=0,\ldots,n-1.
\]
Here $\psi_0$ corresponds to the eigenvalue $0$, while
$\psi_j$, $j=1,\ldots,n-1$, correspond to the eigenvalue $n$. Hence the
eigenvalues of $\Delta_{K_n}^{a,b}=b\Delta_{K_n}-aI$ are
\[
\mu_0=-a,\qquad
\mu_j=bn-a,\quad j=1,\ldots,n-1.
\]
Moreover,
\[
\psi_j(x_0)\overline{\psi_j(x_\ell)}
=
\frac1n e^{-2\pi i j\ell/n}.
\]
Substituting this into \eqref{eq. trace fla} with $x=x_\ell$ and $q=n-2$ yields that
for every $k\in\mathbb N_0$, $a\in\mathbb R$, $b\in\mathbb R\setminus \{0\}$ satisfying \eqref{eq:edge} and $\ell=0,\ldots,n-1$,
\[
\sum_{j=0}^{n-1}
e^{-2\pi i j\ell/n}\mu_j^k=n(b(n-1)-a)^k
\sum_{m=0}^{k}
(-1)^m
b_m(x_\ell)(n-2)^{-m/2}
I_m^{c_{a,b}(n)}(k).
\]
Substituting the coefficients $b_m(x_\ell)$ from \eqref{eq. bm complete graph}
into this identity and simplifying gives that
\begin{multline}\label{eq. ident for kn}
\sum_{j=0}^{n-1}
e^{-2\pi i j\ell/n}
\mu_j^k
=
n(b(n-1)-a)^k
\Bigg[
\delta_{\ell,0} I_0^{c_{a,b}(n)}(k) \\
\quad+
\frac{2}{n}
\sum_{m=1}^{k}
(-1)^m
\left(
T_m\!\left(\frac{n-1}{2\sqrt{n-2}}\right)
+
\eta_\ell
T_m\!\left(-\frac{1}{2\sqrt{n-2}}\right)
\right)
I_m^{c_{a,b}(n)}(k)
\Bigg].
\end{multline}
For $\ell\in\{1,\ldots,n-1\}$, we have $\eta_\ell=-1$. Moreover,
\[
\sum_{j=1}^{n-1}e^{-2\pi ij\ell/n}
=-1.
\]
Since $\mu_0=-a$, $\mu_j=bn-a$ for $j=1,\ldots,n-1$, it follows that
\[
\sum_{j=0}^{n-1}
e^{-2\pi ij\ell/n}\mu_j^k
=
(-a)^k-(bn-a)^k.
\]
Taking $\ell\in\{1,\ldots,n-1\}$ in \eqref{eq. ident for kn}, we obtain
\[
\sum_{m=1}^{k}
(-1)^m
\left(
T_m\!\left(\frac{n-1}{2\sqrt{n-2}}\right)
-
T_m\!\left(-\frac{1}{2\sqrt{n-2}}\right)
\right)
I_m^{c_{a,b}(n)}(k)
=
\frac{(-a)^k-(bn-a)^k}{2(b(n-1)-a)^k}.
\]
Taking $\ell=0$ in \eqref{eq. ident for kn} and observing that $\eta_0=n-1$, and
\[
\sum_{j=0}^{n-1}\mu_j^k
=
(-a)^k+(n-1)(bn-a)^k,
\]
we obtain
\begin{multline*}
\sum_{m=1}^{k}
(-1)^m
\left(
T_m\!\left(\frac{n-1}{2\sqrt{n-2}}\right)
+
(n-1)T_m\!\left(-\frac{1}{2\sqrt{n-2}}\right)
\right)
I_m^{c_{a,b}(n)}(k)\\
=
\frac{(-a)^k+(n-1)(bn-a)^k}{
2(b(n-1)-a)^k}
-\frac{n}{2}I_0^{c_{a,b}(n)}(k).
\end{multline*}
Now put $D=b(n-1)-a$ and $c=\frac{2b\sqrt{n-2}}{D}$. Then
\[
-a
=
D\left(1-\frac{(n-1)c}{2\sqrt{n-2}}\right),
\qquad
bn-a
=
D\left(1+\frac{c}{2\sqrt{n-2}}\right).
\]
Substituting these relations into the last two displayed identities gives the stated formulas valid in the range of $c$ for which $a,\, b$ satisfy \eqref{eq:edge}. This range contains infinitely many $c$, and hence extension to all real $c$ follows from the fact that both sides of equations \eqref{eq:Kn_first_identity} and \eqref{eq:Kn_second_identity}  are polynomials in $c$.
\end{proof}

\end{document}
